\documentclass[11pt]{article}
\usepackage{amsmath,amssymb,amsthm,hyperref}
\newtheorem{theorem}{Theorem}
\newtheorem{lemma}{Lemma}
\newcommand{\E}{\mathbb E}
\title{A polynomial rate for the uniform endpoint profile}
\author{Complete asymptotic proof; independently reviewed}
\date{30 September 2026}
\begin{document}\maketitle
Assume $p_1,\ldots,p_m:[K]\to[0,1]$ are nondecreasing, with positive
probability row weights $w$, and
\[
 \E_w\prod_{j\in S}p_j=\frac1{|S|+1}\qquad(S\subseteq[m]).
\]
Put $u_i=1-p_i$ and $s=\sum_i u_i$.

\begin{theorem}\label{profile}
For every fixed $0<A<1$, there is a constant $C_A$ such that, for
all sufficiently large $m$, uniformly over columns and rows with
$s\le A\log m$,
\[
 \max_i|m(1-p_i)-s|
       \le C_A m^{-(1-A)/8}\log m.                 \tag{1}
\]
In particular this is a quantitative bound on the range of the
comparison columns, rather than only a bound on their maximum.
No new dice-face exponent is asserted in this note.
\end{theorem}

We use the already reviewed marked comparison from
\path{quantitative_relative_endpoint_balance.tex}. For a fixed
column $i$, put $M=m-1$, $\ell=\lfloor M/4\rfloor$,
$\alpha=\ell/M$, $s_-=\sum_{j\ne i}u_j$, and $x=\alpha s_-$.
For a uniform $\ell$-subset $B$ of the remaining columns let
$\overline A=\E_B\prod_{j\in B}p_j$. The two probability measures
\[
 \mu_0=(\ell+1)w\overline A,\qquad
 \mu_1=(\ell+1)(\ell+2)w u_i\overline A
\]
have density relation $d\mu_1=(\ell+2)u_i\,d\mu_0$.
Their scalar targets have densities
\[
 b_0(z)=\frac{\ell+1}{\ell}(1-z/\ell)^\ell\boldsymbol1_{[0,\ell]},
 \qquad b_1(z)=\frac{\ell+2}{\ell}z b_0(z).
 \tag{2}
\]
Denote their integrals by $\beta_0,\beta_1$, and fix $\theta=1/16$.
All constants in the auxiliary bounds below are uniform in $i$.

\begin{lemma}[The degree range in the marked comparison]\label{degree}
For some absolute $c_0>0$, the marked bilinear comparison holds for
all degrees $D\le c_0M$:
\[
 |\mu_a(Q_1Q_2)-\beta_a(Q_1Q_2)|
 \le C B_1B_2(D/M+e^{-cM}),\qquad a=0,1,          \tag{3}
\]
where $Q_b=\sum_{j=0}^D c_{bj}L_j(\theta z)$ and
$B_b=\sum_j|c_{bj}|$. Also
\[
 \mu_a(e^{\theta x/2})+\beta_a(e^{\theta z/2})\le C.
 \tag{4}
\]
\end{lemma}
\begin{proof}
In the proof of the cited comparison, the polarization series is
bounded by a constant times
\[
 B_1B_2\sum_{k=2}^{2D}(C\sqrt{D/M})^k(k+1)^C.
\]
The exponent on $k+1$ is fixed, also for the marked measure. Taking
$c_0$ sufficiently small makes the geometric ratio at most $1/4$,
so this remains $O(B_1B_2D/M)$. Choose also $2c_0M\le M-\ell$.
The weighted-subset comparison and its exponential-tail argument
have no smaller degree restriction; $D\le c_0M$ is polynomially
bounded and its tail remains exponentially small. This proves (3).
Bound (4) follows from $\overline A\le e^{-\alpha s_-}$, the row-mass
Laplace bound, and $u_i\le C(s_-+1)/M$ for the marked normalization.
It is also immediate for the scalar densities (2).
\end{proof}

\section*{An elementary Laguerre smoothing kernel}
Let $0<r<1$, $\tau=1-r$. Given $X=x\ge0$, choose
\[
 N\sim\operatorname{Poisson}(r\theta x/\tau),
 \qquad Y\mid N\sim(\tau/\theta)\operatorname{Gamma}(N+1,1).
 \tag{5}
\]
Write $T_rf(x)=\E[f(Y)\mid X=x]$. The following standard Laguerre
kernel facts are included with proofs. The associated classical
bilinear formula is the Hille--Hardy formula; see the primary
reference \cite[18.18.27]{DLMF}.

\begin{lemma}\label{kernel}
The kernel $T_r$ is stochastically increasing in $x$, reversible for
$\rho_\theta(dx)=\theta e^{-\theta x}dx$, and satisfies
\[
 T_r L_j(\theta x)=r^jL_j(\theta x).
\]
If $f$ is an indicator function and
$a_j=\int f(x)L_j(\theta x)\,d\rho_\theta(x)$, then
$\sum_j a_j^2\le1$ and
\[
 T_rf(x)=\sum_{j=0}^{\infty}r^j a_jL_j(\theta x).
 \tag{6}
\]
After truncation at degree $D$, the remainder is at most
\[
 C\tau^{-1/2}r^{D+1}e^{\theta x/2}.               \tag{7}
\]
\end{lemma}
\begin{proof}
Increasing $x$ can be coupled by adding an independent Poisson
variable to $N$, then adding the corresponding independent unit-rate
exponentials to the gamma sum. This proves stochastic monotonicity.
Summing the gamma densities gives the transition density
\[
 \frac\theta\tau e^{-(\theta y+r\theta x)/\tau}
 \sum_{n=0}^{\infty}\frac{(r\theta^2xy/\tau^2)^n}{(n!)^2}.
\]
Multiplication by $\theta e^{-\theta x}$ makes this symmetric in
$x,y$, proving reversibility and stationarity. The conditional moment
generating function is
\[
 \E(e^{tY}\mid x)
 =(1-\tau t/\theta)^{-1}
       \exp\left(\frac{rxt}{1-\tau t/\theta}\right).
 \tag{8}
\]
Applying this to the ordinary Laguerre generating function
$(1-z)^{-1}\exp[-\theta xz/(1-z)]$ gives the same generating function
with $z$ replaced by $rz$, proving the eigenfunction identity.
The Laguerre polynomials form an orthonormal basis of
$L^2(\rho_\theta)$. Thus reversibility and the eigenfunction identity
prove (6) in $L^2$, with Bessel's inequality giving the coefficient
bound. The series converges absolutely and uniformly on compact
intervals, since $|L_j(\theta x)|\le e^{\theta x/2}$; the transition
integral is continuous for an interval indicator, and hence agrees
pointwise with that series. Alternatively the displayed transition
kernel and its Laguerre expansion give the same identity directly.
Cauchy--Schwarz bounds the tail by
$e^{\theta x/2}(\sum_{j>D}r^{2j})^{1/2}$, proving (7).
Only interval indicators are needed below.
\end{proof}

\begin{lemma}[Uniform crossing concentration]\label{concentration}
There are absolute constants $C,c>0$ such that, if $L\ge1$,
$0<h\le1$, $h\ge C\tau(L+1)$, and $0\le x\le L+1$, then
\[
 \Pr(|Y-x|>h\mid x)\le2\exp[-c h^2/(\tau(L+1))].
 \tag{9}
\]
\end{lemma}
\begin{proof}
Formula (8) gives $\E(Y\mid x)=rx+\tau/\theta$. For
$|t|\tau/\theta\le1/2$, subtracting this mean in its logarithm gives
\[
 \log\E[e^{t(Y-\E Y)}\mid x]
 \le C\tau(x+\tau)t^2.
\]
Indeed $-\log(1-v)-v\le C v^2$ and
$t/(1-v)-t=tv/(1-v)$, with $v=\tau t/\theta$.
The deterministic mean displacement has magnitude at most
$C\tau(L+1)$. A Chernoff bound with
$|t|=c h/(\tau(L+1))$, after increasing the constant in the
hypothesis on $h$, proves (9).
\end{proof}

\section*{Distribution functions with a polynomial error}
\begin{lemma}\label{cdf}
For $a=0,1$ and $1\le L\le\log M$, uniformly in $0\le z\le L$,
\[
 |\mu_a(x\le z)-\beta_a([0,z])|
 \le \varepsilon_L:=C M^{-1/4}\sqrt{(L+1)\log M}.
 \tag{10}
\]
\end{lemma}
\begin{proof}
Take $\tau=M^{-1/2}$ and truncate (6) at
$D=\lfloor c_0M\rfloor$. Applying (3) to the individual products
$1\cdot L_j$, summing, and using Cauchy--Schwarz gives, for any
interval indicator $f$,
\begin{align*}
 |\mu_a(T_rf)-\beta_a(T_rf)|
 &\le\frac C M\sum_{j\le D}j r^j|a_j|
       +C\tau^{-1/2}e^{-c\tau M}\\
 &\le\frac C{M\tau^{3/2}}
       +C\tau^{-1/2}e^{-c\tau M}
 \le C M^{-1/4}.
\end{align*}
The tails use (4) and (7); the constant polynomial has zero
comparison error. The series $\sum j^2r^{2j}$ is $O(\tau^{-3})$.

Choose $h=C_1\sqrt{\tau(L+2)\log M}$ with $C_1$ sufficiently large.
For large $M$, it is at most $1/4$ and satisfies the hypothesis of
Lemma~\ref{concentration}. The crossing error is at most $M^{-10}$.
For each threshold $z\in[0,L]$, stochastic monotonicity and (9)
therefore imply
\[
 (1-M^{-10})\boldsymbol1_{x\le z}
 \le T_r\boldsymbol1_{[0,z+h]}(x)
 \le\boldsymbol1_{x\le z+2h}+M^{-10}.
\]
For $x$ outside the concentration interval, stochastic monotonicity
reduces the bound to $x=z$ or $x=z+2h$; thus these inequalities hold
on the whole half-line. The reverse comparison follows by shifting
the threshold down by $h$, with a trivial lower bound when $z<2h$.
Integrate these inequalities against $\mu_a$ and $\beta_a$, use the
smoothed-test comparison above, and note that both scalar densities
in (2) are bounded by an absolute constant. This proves (10).
It also handles atoms of $\mu_a$ and either choice of strict endpoints.
\end{proof}

\section*{Recovering the monotone density}
\begin{lemma}\label{ratio}
Suppose $1\le L\le\log M$ and
\[
 h_*=\sqrt{(L+1)\varepsilon_{L+2}e^L}\longrightarrow0
\]
along the range under consideration. Uniformly over rows with $x\le L$,
\[
 |\ell u_i-x|\le C h_* .                         \tag{11}
\]
Here (10) can be used with $L+2$; replacing the upper limit
$\log M$ there by $\log M+2$ only changes constants.
\end{lemma}
\begin{proof}
For $0\le z\le L+2$, the explicit scalar density satisfies
$b_0(z)\ge c e^{-z}$ for large $M$, because
$L=O(\log M)$ and
$\ell\log(1-z/\ell)=-z+O((\log M)^2/M)$.
Choose an interval of length $h_*$ lying strictly above the target
row's value $x$, at distance between $h_*$ and $2h_*$.
Every row in it lies before the target row, so monotonicity and the
exact density relation give
\[
 (\ell+2)u_i(q)\le\frac{\mu_1(I)}{\mu_0(I)}.
\]
The scalar denominator is at least $c h_*e^{-L}$, and the two
interval masses have errors at most $2\varepsilon_{L+2}$ by (10).
For large $M$ the denominator error is less than half the scalar
mass, since
$\varepsilon_{L+2}e^L/h_*=h_* /(L+1)\to0$.
The scalar density ratio $b_1/b_0=((\ell+2)/\ell)z$ varies by
$O(h_*)$ across the chosen interval. Perturbing numerator and
denominator therefore bounds its ratio error by
\[
 C(L+1)\varepsilon_{L+2}e^L/h_*=C h_*.
\]
This proves the upper bound in (11). For $x\ge2h_*$ use the
corresponding interval strictly below $x$ for the lower bound.
When $x<2h_*$, nonnegativity supplies that bound directly.
Flat row blocks cause no problem because the comparison intervals
are strictly separated from the target value.
\end{proof}

\begin{proof}[Proof of Theorem~\ref{profile}]
For a row with $s\le A\log m$ one has
$x=\alpha s_-\le(A/4)\log m$. Set
$L=\max(1,(A/4)\log m)$. Then
\[
 h_*\le C_A M^{-1/8}e^{L/2}(L+1)^{3/4}
                                      (\log M)^{1/4}
       \le C_A m^{-(1-A)/8}\log m\longrightarrow0.
\]
Lemma~\ref{ratio} applies uniformly in the deleted index $i$.
Finally the conversion to the full mean is exact:
\[
 \frac M\ell(\ell u_i-x)=M u_i-s_-=m u_i-s.
\]
Since $M/\ell$ is bounded, this gives (1).
\end{proof}

\begin{thebibliography}{1}
\bibitem{DLMF} NIST Digital Library of Mathematical Functions,
\emph{Hille--Hardy formula}, equation 18.18.27.
\url{https://dlmf.nist.gov/18.18.E27}.
\end{thebibliography}
\end{document}
